\scanpage{160}
Groupe fondamental dans les catégories.

$\unsure{\mathcal C|\mathscr C}$ une catégorie, avec une famille $(D)$ de morphismes
de $\mathcal C$, satisfaisant les conditions usuelles

(i) contient identité et stable par composition (ii) stable par
changement de base

On suppose de plus

(iii) \deleted{\unsure{les}} \inserted{$D$} sont des morphismes de descente [donc des
\unsure{épim.} stricts \inserted{universels} si \inserted{\unsure{on suppose que}} \uline{les produits et produits}
\uline{fibrés existent}]. Enfin on suppose

% slanted
\margin{On peut \unsure{axer|varier|aussi} la situation, en introduisant de plus \deleted{\unsure{la}} une classe de morphismes $D'$ [\deleted{de} morphismes pouvant être \unsure{``localement''} \unsure{fibrés}.}

(iv) On peut ``descendre'' un $g\in(D)$ par un $f\in(D)$
[mais on ne suppose pas qu'on peut descendre
\uline{tout} morphisme par un $f\in(D)$]. \uline{Dans la
suite, $D$ est supposé donné une fois pour toutes.}
\uline{On suppose que dans $\mathcal C$ les sommes finies} [\unsure{resp.|évent.} quelconques] existent\unsure{]|.}

\margin{(v) La projection $X\times I\to X$ est $\in(D)$.}

\uline{Exemple} Catégorie \inserted{et \unsure{première?} condition} $\unsure{\mathscr E|\mathcal E|\mathcal C}$ des préschémas. \deleted{La catégorie}
Classe $(D)$ des morphismes \uline{plats} et \uline{affines}.
[ou \uline{plats} et quasi-projectifs].

On va développer une théorie particulière des
fibrés principaux [vérifiant les conditions générales
\deleted{\unsure{Groupes de base|Groupes de faisceaux}}
d'une théorie des fibrés principaux, cf. \unsure{---}].

\uline{Groupes structuraux} : $X\times G$ $G$ un groupe ensembliste
fini [\unsure{resp.} quelconque]

\uline{Fibrés principaux \unsure{locaux} sur $X$}, de groupe $G$ : les
fibrés \uline{\unsure{fonctoriels|formels|fonctls|formés}} ainsi, avec hom $P\mathrel{\unsure{-|\to}}X$ $\in(D)$
% archive sheet number: 159
